\part*{Part A}
\addcontentsline{toc}{part}{Part A}

\noindent\textbf{Chosen Topic for Part A:} Sleep quality of Singapore residents: findings from the 2016 Singapore mental health study \cite{lee2022}.

% ============================================================
\section*{Section 1}
\addcontentsline{toc}{section}{Section 1}

\subsection*{1. Aim of Study}
\addcontentsline{toc}{subsection}{1. Aim of Study}

The primary objective of this study was to investigate the quality of sleep among community dwellers in Singapore. Specifically, the researchers wanted to:
\begin{enumerate}
    \item Characterise the sleep quality of adult Singaporean residents in a nationally representative sample.
    \item Identify the sociodemographic factors that correlate with poor sleep in this population. The sociodemographic correlates in the study include Gender, Ethnicity, Marital status, Education, Lifestyle (smoking, drinking, etc.), Mental health illnesses, and physical health illnesses.
\end{enumerate}

\subsection*{2. Main Independent and Dependent Variables}
\addcontentsline{toc}{subsection}{2. Main Independent and Dependent Variables}

\noindent\textbf{Dependent Variables:}

\noindent\textit{Sleep Quality:} Measured using the Pittsburgh Sleep Quality Index (PSQI), a self-report questionnaire that assesses sleep patterns over a one-month interval.

The PSQI is a standardized clinical tool that asks about seven components: subjective sleep quality, sleep latency (time to fall asleep), sleep duration, habitual sleep efficiency, sleep disturbances, use of sleeping medication, and daytime functioning.

A global score was calculated from 0--21; a total score of 5 or more was used to define ``poor sleepers,'' while a score below 5 indicated ``good sleepers.'' Sleep duration obtained was then classified into $\leq 6$h, 7--8h and $\geq 9$h.

\noindent\textbf{Independent Variables:}
\begin{itemize}
    \item \textit{Mental Disorders:} Measured using the World Health Organization Composite International Diagnostic Interview (WHO CIDI 3.0) to diagnose conditions such as Major Depressive Disorder and Obsessive-Compulsive Disorder. Respondents were categorized into three types: no mental illness, at least one mental illness, and two or more comorbid mental illnesses in the past 12 months.

    \item \textit{Physical Disorders:} Measured by participant self-report of chronic medical conditions (e.g.\ diabetes, heart disease) that had been diagnosed by a doctor. Responses to the checklist were grouped into 3 types: no chronic physical conditions, at least one physical chronic condition, and multimorbidity in the past 12 months. Participants were given a list of 18 common chronic medical conditions in Singapore such as hypertension, diabetes, asthma, and neurological disorders and were asked if a healthcare professional had ever told them they had the condition.
\end{itemize}

The researchers state that the observed stronger and more consistent associations between sleep and the above two independent variables. Hence, we believe that mental and physical disorders were the main 2 independent variables, while the other sociodemographic factors were considered secondary independent variables.

\subsection*{3. Main Findings of the Article}
\addcontentsline{toc}{subsection}{3. Main Findings of the Article}

\begin{enumerate}[label=\alph*)]
    \item ``In this sample, only certain measures of socioeconomic status were moderately associated with sleep quality and duration. Instead, we observed stronger and more consistent associations between sleep and physical and mental health conditions" \cite{lee2022}.

    \item The target population were residents of Singapore who were adults (including citizens and permanent residents) aged 18 years and above. The study used disproportionate stratified sampling across three main ethnic groups (Chinese, Malay, Indian) and oversampled those aged 65 and above, with sampling weights then applied to make the findings generalisable to the broader adult Singapore population.

    \item To examine whether mental health conditions are significantly associated with poor sleep quality among Singapore residents, the following hypotheses were formulated:

    \noindent$H_0$: There is no significant association between mental health conditions and poor sleep quality among Singapore residents.

    \noindent$H_1$: There is a significant association between mental health conditions and poor sleep quality among Singapore residents.

    The results of the chi-square, logistic regression and multimodal analysis tests revealed a statistically significant association between mental health conditions and poor sleep quality. Relative to individuals with no mental illness, those diagnosed with one mental illness demonstrated significantly higher odds of reporting poor sleep quality (AOR $= 3.42$, 95\% CI $[2.36,\, 4.95]$, $p < 0.001$). Furthermore, individuals with comorbid mental illnesses exhibited a substantially greater likelihood of poor sleep quality (AOR $= 14.11$, 95\% CI $[6.52,\, 30.54]$, $p < 0.001$), suggesting a dose-dependent relationship between the number of mental health conditions and the severity of sleep impairment.

    \noindent\textbf{Decision:} As both $p$-values were well below the predetermined significance level of $\alpha = 0.05$, and neither confidence interval contained the value of 1, there was sufficient statistical evidence to reject the null hypothesis. The findings indicate that mental health conditions are significantly and positively associated with poor sleep quality among Singapore residents.
\end{enumerate}

% ============================================================
\section*{Section 2}
\addcontentsline{toc}{section}{Section 2}

\subsection*{4. Confounders in the Observational Study}
\addcontentsline{toc}{subsection}{4. Confounders in the Observational Study}

This was an observational study that aimed to identify socio-demographic factors correlated with sleep quality. The researchers also included clinical health issues and smoking and drinking habits as areas to look into. These qualities are considered independent variables for the study. There were no specific confounders controlled by the researchers, but these independent variables could serve as confounders for each other. The confounding independent variables that are controlled by the researchers include mental illnesses, chronic physical illnesses, ethnicity, gender, education level, household income, age, smoking habits, drinking habits, marital status, employment status and body mass index.

Gender is a confounding variable. Being female was associated with both higher occurrence of mental illness and lower sleep quality, thus gender may confound the association between mental illness and sleep quality.

A potential confounder not controlled by researchers is whether the participant's house is located near sources of nighttime noise, such as busy roads or areas with active nightlife. Loud noises may disrupt sleep, leading to poor sleep quality.

\subsection*{5. Selection Process}
\addcontentsline{toc}{subsection}{5. Selection Process}

The subjects were Singapore citizens and permanent residents over the age of 18. There were 6,126 subjects with a response rate of 69.5\%. This sampling frame covers the general population of Singapore residents fairly well. However, it excludes children younger than 18 years old, residents unable to attend the face-to-face interview due to being out of the country, being institutionalized or hospitalized, residents unable to complete the interview due to physical and/or mental illnesses, language barriers, and residents that could not be contacted using information from the national registry. Probability sampling was employed in the selection of respondents.

Due to its failure to account for certain groups of the target population, the study may lack crucial information from these non-responding or excluded subjects. Some groups of the population may have a rarer occurrence, such as chronic physical illness. With a sample size of 6,126, this may cause these smaller subgroups to have less or even no representation in the study, leading to insufficient evidence for reliable results in these groups. Thus, this study cannot be fully generalized to the target population of Singapore residents.

\subsection*{6. Concluding Evaluation}
\addcontentsline{toc}{subsection}{6. Concluding Evaluation}

This study put in good efforts in conducting their study well. The researchers did their best to minimise non-response through repeated attempts at contacting selected respondents at different timings. The study also uses disproportionate stratified sampling to select more subjects from the elderly and minority ethnicities, then applies sampling weights to adjust the results to represent Singapore's population. As such, this allowed enough data of these groups with lesser numbers to be collected for more reliable results.

However, the study had some weaknesses. It relied on participants to self-report their quality of sleep through the Pittsburgh Sleep Quality Index. While this method was chosen for its feasibility, participants may misrecall their sleep characteristics or be affected by recency or other biases, causing inaccurate reporting of their sleep quality and some unreliability. The data were also collected from 2016 to 2018, making them relatively old compared to the current 2026 context.

The study is generally well-conducted. The strengths in having a relevant sampling frame, random sampling, and stratified sampling of some minority groups outweigh the weaknesses of the study method, older data and lack of generalisability for smaller subgroups. By having a sample representative of the general population rather than a specific demographic, it provides insights into the sleep quality of Singapore residents.